\newpage
\section*{Solucionario completo}

\noindent Cada respuesta incluye la clave y una justificación breve. En los ejercicios numéricos se indica la ecuación o el cálculo principal.

\subsection*{Bloque A. Fundamentos teóricos y conceptuales}

\sol{1}{D}{El equilibrio general considera simultáneamente todos los mercados de la economía. Los precios se ajustan hasta que la oferta y la demanda coinciden en cada mercado.}

\sol{2}{B}{Un equilibrio walrasiano es un vector de precios tal que todos los mercados se vacían. No es cualquier asignación eficiente ni un precio fijado por el Estado.}

\sol{3}{C}{Las funciones de demanda y oferta dependen de los precios relativos. Además, se puede elegir un bien como numerario, de modo que solo importan $k-1$ precios independientes.}

\sol{4}{A}{La ley de Walras indica que la suma del valor de los excesos de demanda es cero: $\sum_i p_i z_i(p)=0$. Si $n-1$ mercados se vacían, el mercado restante también lo hace.}

\sol{5}{D}{La función de exceso de demanda es homogénea de grado cero en precios: si todos los precios se multiplican por una constante positiva, las cantidades demandadas y ofrecidas no cambian.}

\sol{6}{B}{La existencia de equilibrio requiere continuidad de las funciones de exceso de demanda, convexidad de las preferencias y otros supuestos estándar. El monopolio o los precios rígidos no forman parte de esos supuestos.}

\sol{7}{C}{El equilibrio parcial analiza un mercado aislando los demás; el equilibrio general incorpora las interacciones y los efectos de retroalimentación entre mercados.}

\sol{8}{D}{Una asignación es Pareto eficiente cuando no existe otra asignación factible que mejore a al menos una persona sin empeorar a otra.}

\sol{9}{D}{Una mejora de Pareto exige que al menos un individuo esté mejor y que ningún otro esté peor.}

\sol{10}{A}{El Primer Teorema del Bienestar señala que un equilibrio competitivo es Pareto eficiente bajo supuestos como mercados completos, competencia perfecta y ausencia de externalidades.}

\sol{11}{B}{Las externalidades, los bienes públicos, el poder de mercado y la información asimétrica son fallas de mercado que pueden invalidar el Primer Teorema del Bienestar.}

\sol{12}{C}{El Segundo Teorema del Bienestar establece que, bajo supuestos de convexidad y mercados completos, cualquier asignación Pareto eficiente puede descentralizarse como un equilibrio competitivo con transferencias lump-sum.}

\sol{13}{A}{Las transferencias lump-sum no dependen de las decisiones de los agentes y no crean distorsiones marginales; solo redistribuyen ingreso entre ellos.}

\sol{14}{B}{La caja de Edgeworth representa todas las asignaciones factibles de dos bienes entre dos consumidores, dadas las dotaciones totales.}

\sol{15}{B}{El punto de dotación es la asignación inicial de recursos de la que parten los consumidores antes de intercambiar.}

\sol{16}{A}{La curva de contrato es el conjunto de asignaciones Pareto eficientes. En intercambio puro, corresponde a los puntos donde las curvas de indiferencia son tangentes.}

\sol{17}{A}{El núcleo o core está formado por las asignaciones que ninguna coalición de agentes puede mejorar mediante sus propios recursos.}

\sol{18}{D}{La lente está delimitada por las curvas de indiferencia que pasan por la dotación inicial; todos sus puntos mejoran a ambos consumidores respecto de esa dotación.}

\sol{19}{D}{En un óptimo interior de intercambio puro, las curvas de indiferencia son tangentes y, por tanto, $MRS_A=MRS_B$.}

\sol{20}{B}{Con producción, la eficiencia requiere que la tasa marginal de sustitución entre bienes sea igual a la tasa marginal de transformación: $MRS=MRT$.}

\sol{21}{C}{La frontera de posibilidades de utilidad muestra las combinaciones máximas de utilidad de dos individuos; indica la máxima utilidad de uno dada la utilidad del otro.}

\sol{22}{B}{Una función de bienestar social agrega las utilidades individuales en un ordenamiento social. La elección de la función implica un juicio distributivo.}

\sol{23}{A}{La eficiencia de Pareto se refiere a la imposibilidad de mejorar sin empeorar a otro; no dice nada sobre la igualdad o la justicia de la distribución.}

\sol{24}{C}{Una transferencia lump-sum redistribuye ingreso sin modificar los incentivos marginales ni introducir distorsiones en las decisiones de consumo o producción.}

\sol{25}{D}{El criterio de Kaldor-Hicks acepta una política si los ganadores podrían compensar a los perdedores y aun así mejorar, aunque la compensación no se pague efectivamente.}

\sol{26}{A}{Una externalidad negativa hace que el costo social marginal difiera del costo privado; el equilibrio competitivo no internaliza ese costo y no es Pareto eficiente.}

\sol{27}{A}{Un bien público puro no es rival ni excluible. El mercado tiende a subproveerlo porque los individuos pueden beneficiarse sin pagar.}

\sol{28}{C}{El monopolio fija $P>CMg$ y produce menos que la cantidad competitiva, lo que genera una pérdida de bienestar.}

\sol{29}{A}{El Primer Teorema del Bienestar requiere mercados completos, competencia perfecta, información perfecta, ausencia de externalidades y otros supuestos.}

\sol{30}{C}{Con mercados incompletos o información asimétrica, el equilibrio competitivo puede no ser Pareto eficiente; el Primer Teorema deja de aplicarse automáticamente.}

\subsection*{Bloque B. Aplicaciones cuantitativas y optimización}

\sol{31}{C}{La caja de Edgeworth tiene ancho igual a la dotación total del bien $x$: $X=10$.}

\sol{32}{D}{La altura de la caja de Edgeworth es la dotación total del bien $y$: $Y=10$.}

\sol{33}{B}{La dotación inicial de A es la pareja $(8,2)$. La de B es $(2,8)$, que completa las dotaciones totales.}

\sol{34}{D}{Para $U_A=x_Ay_A$, $MRS_A=MU_x/MU_y=y_A/x_A$. En la dotación: $2/8=0.25$.}

\sol{35}{D}{Para B, $MRS_B=y_B/x_B=8/2=4$.}

\sol{36}{B}{En la dotación, $MRS_A=0.25$ y $MRS_B=4$. Como son distintas, existen ganancias de intercambio y la asignación no es Pareto eficiente.}

\sol{37}{D}{La eficiencia exige $y_A/x_A=y_B/x_B$. Con $x_B=10-x_A$ e $y_B=10-y_A$ se obtiene $y_A=x_A$.}

\sol{38}{B}{Con preferencias Cobb-Douglas idénticas y dotaciones simétricas, el precio relativo de equilibrio es $p_x/p_y=1$.}

\sol{39}{C}{Con $p_x=p_y=1$, el ingreso de A es $8+2=10$. Gasta la mitad en $x$: $x_A=5$.}

\sol{40}{C}{Con el mismo ingreso y precios, A gasta la otra mitad en $y$: $y_A=5$. El equilibrio de A es $(5,5)$.}

\sol{41}{D}{En la dotación, $U_A=8\cdot 2=16$. En el equilibrio, $U_A=5\cdot 5=25$. La ganancia es $25-16=9$.}

\sol{42}{C}{Para B, $U_B=2\cdot 8=16$ en la dotación y $U_B=5\cdot 5=25$ en el equilibrio; la ganancia también es 9.}

\sol{43}{C}{En el equilibrio, A consume $(5,5)$ y B consume $(5,5)$. Entonces $MRS_A=5/5=1$ y $MRS_B=5/5=1$; la asignación es Pareto eficiente.}

\sol{44}{A}{Con $p_x=p_y=1$, la asignación A$=(6,6)$ requiere un ingreso de 12. A tiene un ingreso inicial de 10, por lo que necesita una transferencia de 2 desde B.}

\sol{45}{D}{Con dotación A$=(9,1)$ y B$=(1,9)$, a precios unitarios ambos tienen ingreso 10. A consume $x_A=5$ e $y_A=5$.}

\sol{46}{C}{Con una transferencia de 2 desde B hacia A, el ingreso de A es 12 y el de B es 8. A consume $x_A=6$ e $y_A=6$.}

\sol{47}{A}{La asignación A$=(6,6)$ y B$=(4,4)$ cumple $MRS_A=6/6=1$ y $MRS_B=4/4=1$. Las MRS son iguales, por lo que es Pareto eficiente.}

\sol{48}{A}{El Segundo Teorema del Bienestar permite sostener una asignación Pareto eficiente como equilibrio competitivo mediante transferencias lump-sum, bajo supuestos de convexidad y mercados completos.}

\sol{49}{D}{Si el total de $x$ es 20 y A tiene 12, entonces B tiene $20-12=8$.}

\sol{50}{B}{Si el total de $y$ es 10 y A tiene 4, entonces B tiene $10-4=6$.}

\sol{51}{A}{Con preferencias convexas, en un óptimo interior de intercambio la igualdad de las MRS es la condición de tangencia y, por tanto, de eficiencia de Pareto.}

\sol{52}{D}{Si $MRS_A=3$ y $MRS_B=1$, las valoraciones marginales difieren; existe un intercambio mutuamente beneficioso hasta igualar las MRS.}

\sol{53}{A}{La ley de Walras con dos bienes y precios unitarios implica $z_x+z_y=0$. Si $z_x=4$, entonces $z_y=-4$.}

\sol{54}{C}{Con $p_x=p_y=1$ y $z_x=12$, la ley de Walras implica $z_y=-12$.}

\sol{55}{C}{Si $p_x/p_y=2$ y $p_y=1$, entonces $p_x=2$.}

\sol{56}{B}{El precio relativo es $p_x/p_y=3/6=0.5$.}

\sol{57}{C}{Si $U_A+U_B=100$ y $U_B\geq 0$, el máximo de $U_A$ se alcanza cuando $U_B=0$: $U_A=100$.}

\sol{58}{C}{Para $W=\min(U_A,U_B)$ con $U_A+U_B=100$, el máximo se alcanza en el punto igualitario: $U_A=U_B=50$.}

\sol{59}{D}{Para $W=U_A+U_B$ y $U_A+U_B=100$, cualquier punto de la frontera da $W=100$. La distribución no queda determinada por el criterio utilitarista.}

\sol{60}{B}{La función rawlsiana maximiza la utilidad del individuo peor situado; se concentra en el mínimo de las utilidades.}

\sol{61}{A}{El Primer Teorema del Bienestar requiere mercados completos, ausencia de externalidades y agentes precio-aceptantes, entre otros supuestos.}

\sol{62}{B}{Con una externalidad negativa, el costo social marginal no coincide con el privado; el equilibrio competitivo no es Pareto eficiente.}

\sol{63}{A}{El impuesto pigouviano óptimo iguala el costo externo marginal, de modo que el agente internaliza el daño que genera.}

\sol{64}{C}{La condición de Samuelson para un bien público es que la suma de las tasas marginales de sustitución iguale la tasa marginal de transformación: $\sum MRS=MRT$.}

\sol{65}{C}{En un monopolio, $P>CMg$ y la cantidad producida es menor que la competitiva; esto genera pérdida de bienestar y una falla del Primer Teorema del Bienestar.}

\subsection*{Bloque C. Casos con implicancias regulatorias y de política}

\sol{66}{B}{La eficiencia requiere $MBA=MBB$: $100-q_A=80-q_B$. Con $q_A+q_B=100$, se obtiene $100-q_A=80-(100-q_A)$, de donde $q_A=60$.}

\sol{67}{B}{Con $q_A=60$ y un total de 100, el usuario B recibe $q_B=40$.}

\sol{68}{C}{El beneficio total inicial es $4550+1950=6500$. En la asignación eficiente es $4200+2400=6600$. La ganancia es 100.}

\sol{69}{B}{Si el total del otro bien es 200 y A tiene 80, B tiene $200-80=120$.}

\sol{70}{A}{Con $MRS_A=3$ y $MRS_B=1$, las valoraciones marginales difieren; existe un intercambio que mejora a ambos.}

\sol{71}{B}{En un óptimo interior de asignación de agua, la tangencia de las curvas de indiferencia implica $MRS_A=MRS_B$.}

\sol{72}{D}{La curva de contrato está formada precisamente por las asignaciones Pareto eficientes.}

\sol{73}{A}{La diferencia entre la asignación inicial A$=80$ y la eficiente A$=60$ es 20 unidades, que se transfieren de A hacia B.}

\sol{74}{A}{La reasignación aumenta el beneficio total, pero A podría perder si no recibe compensación. Existen ganancias potenciales que permiten una mejora de Pareto con compensación.}

\sol{75}{D}{Con derechos bien definidos y costos de transacción nulos, el teorema de Coase sostiene que el intercambio conduce a una asignación eficiente, independientemente de la asignación inicial de derechos.}

\sol{76}{B}{El Primer Teorema del Bienestar requiere ausencia de externalidades, mercados completos y agentes precio-aceptantes, entre otros supuestos.}

\sol{77}{C}{La externalidad negativa hace que el costo social marginal supere al privado; el equilibrio competitivo no internaliza el daño y no es Pareto eficiente.}

\sol{78}{D}{El impuesto pigouviano debe igualar el costo externo marginal para corregir la externalidad.}

\sol{79}{B}{El Segundo Teorema del Bienestar permite sostener una asignación eficiente de agua mediante precios y transferencias lump-sum, bajo ciertos supuestos.}

\sol{80}{D}{Una transferencia lump-sum focalizada protege el ingreso de los hogares vulnerables sin distorsionar los precios relativos ni generar escasez.}

\sol{81}{D}{En un monopolio natural, $P=CMg<CMe$ genera pérdidas; se requiere un subsidio, una tarifa en dos partes u otro mecanismo de cobertura.}

\sol{82}{C}{La protección de una cuenca es un bien público si no es rival ni excluible; el mercado tiende a subproveerla porque hay incentivos a no pagar.}

\sol{83}{A}{La condición de Samuelson para un bien público es $\sum MRS=MRT$; la suma de las disposiciones marginales a pagar debe igualar el costo marginal de provisión.}

\sol{84}{B}{El criterio de Kaldor-Hicks evalúa una política como deseable si los ganadores podrían compensar a los perdedores y aun así mejorar, aunque la compensación no se realice.}

\sol{85}{A}{Si la compensación se paga efectivamente y nadie queda peor, el resultado es una mejora de Pareto, no solo una mejora potencial.}

\sol{86}{B}{La eficiencia de Pareto no garantiza equidad. Una asignación puede ser eficiente y, al mismo tiempo, muy desigual; se requiere un juicio distributivo adicional.}

\sol{87}{C}{La función rawlsiana maximiza la utilidad del individuo peor situado, es decir, se concentra en el mínimo de las utilidades.}

\sol{88}{A}{La función utilitarista maximiza la suma de las utilidades individuales, sin atender directamente a su distribución.}

\sol{89}{D}{Una reforma tarifaria puede afectar los precios de otros bienes y los ingresos de los hogares; el análisis de equilibrio general incorpora esos efectos en alimentos, energía y trabajo.}

\sol{90}{B}{La evaluación de bienestar debe incluir externalidades y efectos de equilibrio general, además de los excedentes del consumidor y del productor.}

\subsection*{Bloque D. Ampliación: Teoría y Aplicaciones Prácticas (91 a 180)}

\sol{91}{B}{El equilibrio parcial mantiene ceteris paribus los precios de otros bienes y factores; una reforma amplia altera salarios, rentas y precios relativos en múltiples mercados simultáneamente.}
\sol{92}{C}{Por homogeneidad de grado cero ($z(\lambda p) = z(p)$), solo importan los $k-1$ precios relativos; fijar $p_k = 1$ elimina la indeterminación del nivel absoluto de precios.}
\sol{93}{A}{Bajo no saciedad local, cada consumidor gasta todo su ingreso ($p \cdot x_h = p \cdot \omega_h$). Sumando sobre todos los agentes resulta $p \cdot z(p) \equiv 0$ para cualquier $p > 0$.}
\sol{94}{D}{Como $\sum_{i=1}^k p_i z_i(p) = 0$ y $z_1 = \dots = z_{k-1} = 0$, queda $p_k z_k(p) = 0$. Dado $p_k > 0$, necesariamente $z_k(p) = 0$.}
\sol{95}{B}{La Ley de Walras deriva de las restricciones presupuestarias individuales, por lo que vale para todo vector de precios $p > 0$: un exceso de demanda en un mercado implica exceso de oferta en valor en otros.}
\sol{96}{C}{Por holgura complementaria ($p_j^* z_j(p^*) = 0$ con $z_j \le 0$ y $p_j^* \ge 0$), cualquier bien con exceso de oferta estricto en equilibrio tiene precio nulo ($p_j^* = 0$).}
\sol{97}{A}{En el tâtonnement, $\dot{p}_i = \alpha_i z_i(p)$ y no se ejecutan transacciones hasta que el subastador anuncia el vector de precios que vacía todos los mercados.}
\sol{98}{D}{Multiplicar todos los precios por $\lambda = 5$ escala el valor de la dotación y el costo de las canastas por igual, dejando invariantes el conjunto presupuestario y el equilibrio real.}
\sol{99}{B}{Cada punto de la caja representa una asignación factible donde $x_{1A} + x_{1B} = X$ (ancho) y $x_{2A} + x_{2B} = Y$ (altura).}
\sol{100}{C}{Como $A$ valora más el bien $1$ en términos del bien $2$ que $B$ ($TMS^A > TMS^B$), ambos ganan si $A$ adquiere bien $1$ de $B$ y le entrega bien $2$ a una tasa intermedia entre $TMS^B$ y $TMS^A$.}
\sol{101}{A}{En el borde de la caja (por ejemplo, $O_A$, $O_B$ o con sustitutos perfectos), las curvas de indiferencia pueden no ser tangentes ($TMS^A \neq TMS^B$) y aun así ser Pareto eficientes por restricción de no negatividad.}
\sol{102}{D}{La Curva de Contrato reúne todos los Óptimos de Pareto de la caja; el Núcleo restringe ese conjunto a las asignaciones individualmente racionales ($U_A \ge U_A(W)$ y $U_B \ge U_B(W)$).}
\sol{103}{B}{Si una coalición $S$ pudiera mejorar a todos sus miembros con su dotación $\sum_{i\in S}\omega_i$, el costo a precios $p^*$ de esa canasta superaría $\sum_{i\in S} p^*\omega_i$, lo cual es imposible.}
\sol{104}{C}{Al crecer el número de agentes, las coaliciones pueden bloquear cada vez más asignaciones no competitivas; en el límite $r \to \infty$, solo sobreviven los equilibrios walrasianos.}
\sol{105}{A}{Cuando $X > Y$, el excedente $X - Y$ no aporta utilidad adicional a nadie; cualquier reparto de ese sobrante manteniendo $x_A \ge y_A$ y $x_B \ge y_B$ es Pareto eficiente (banda gruesa).}
\sol{106}{D}{La condición $TMS^A = TMS^B$ se reduce a $v_A'(x_{1A}) = v_B'(X - x_{1A})$, que determina un valor único $x_{1A}^*$ independiente de $x_{2A}$ (tramo interior vertical).}
\sol{107}{B}{La curva de oferta de $A$ une sus elecciones óptimas al rotar la recta presupuestaria alrededor de $W$. El equilibrio walrasiano ocurre en la intersección de las curvas de oferta de $A$ y $B$.}
\sol{108}{C}{Maximizar $X$ dado un nivel de $Y$ requiere tangencia entre isocuantas: $TMST_{L,K}^X = TMST_{L,K}^Y$, lo cual garantiza el mercado competitivo al igualar ambas a $w/r$.}
\sol{109}{A}{Mover recursos desde $Y$ hacia $X$ cuesta $dC = CMg_X dX + CMg_Y dY = 0$ sobre la frontera eficiente, de donde $TMT_{X,Y} = -dY/dX = CMg_X / CMg_Y$.}
\sol{110}{D}{Con intensidades factoriales distintas, reasignar factores hacia $X$ libera una proporción inadecuada de $K$ y $L$ respecto de lo que requiere $X$, generando costos de oportunidad crecientes (FPP cóncava).}
\sol{111}{B}{Si las proporciones factoriales óptimas coinciden en ambas industrias y hay RCE (o hay un solo factor), el costo de oportunidad $CMg_X/CMg_Y$ es constante y la FPP es lineal.}
\sol{112}{C}{$TMS_{X,Y} > TMT_{X,Y}$ indica que los consumidores valoran una unidad extra de $X$ en más unidades de $Y$ de las que cuesta producirla; aumentar $X$ a costa de $Y$ genera una mejora de Pareto.}
\sol{113}{A}{Al enfrentar todos los agentes los mismos precios $(p_X, p_Y, w, r)$, se concatenan $TMS^A = TMS^B = p_X/p_Y = CMg_X/CMg_Y = TMT$ y $TMST^X = TMST^Y = w/r$.}
\sol{114}{D}{La FPP entre ocio y consumo es $c = f(T - l)$, cuya pendiente en valor absoluto es $TMT_{l,c} = f'(L) = PMg_L$. En el óptimo, $TMS_{l,c} = PMg_L$ (que en descentralización iguala al salario real $w/p$).}
\sol{115}{B}{Con rendimientos crecientes, $CMg < CMe$, de modo que la regla competitiva $P = CMg$ genera pérdidas y no existe un equilibrio walrasiano precio-aceptante.}
\sol{116}{C}{Todos los consumidores enfrentan $p_X^c = (1+\tau)p_X^p$ (mantienen $TMS^A=TMS^B$) y la producción sigue sobre la FPP, pero $TMS = (1+\tau)TMT \neq TMT$.}
\sol{117}{A}{El Primer Teorema demuestra que todo equilibrio walrasiano es Pareto eficiente usando solo no saciedad local (y ausencia de externalidades/poder de mercado); la convexidad solo se exige en el Segundo Teorema.}
\sol{118}{D}{Un extremo donde un solo agente posee todos los bienes es Pareto eficiente, lo cual ilustra que eficiencia paretiana no implica justicia ni equidad distributiva.}
\sol{119}{B}{La convexidad permite separar el conjunto mejor-que del conjunto factible mediante un hiperplano de precios $p^*$; sin convexidad, el óptimo deseado puede no ser un máximo presupuestario a ningún precio.}
\sol{120}{C}{Solo las transferencias de suma alzada modifican la riqueza inicial sin introducir cuñas entre $TMS$ y $TMT$, preservando todas las condiciones de eficiencia.}
\sol{121}{A}{Al haber asimetría de información sobre la capacidad productiva individual, cualquier impuesto basado en ingresos o consumo observables deja de ser de suma alzada y crea distorsiones (trade-off eficiencia-equidad).}
\sol{122}{D}{El Primer Teorema falla cuando existen fallas de mercado (externalidades, bienes públicos, monopolio, información asimétrica o mercados incompletos), pero no por desigualdad inicial ni preferencias heterogéneas.}
\sol{123}{B}{Cuando una condición paretiana no puede cumplirse, el óptimo restringido (second-best) suele requerir desviarse deliberadamente de las demás condiciones marginales de eficiencia.}
\sol{124}{C}{En teoría, el planificador elige el punto deseado sobre la FPU (equidad) vía transferencias de suma alzada y delega al mecanismo de precios la asignación eficiente.}
\sol{125}{A}{Con mercados completos de bienes contingentes, los agentes comparten eficientemente el riesgo igualando sus tasas marginales de sustitución entre estados a los precios relativos contingentes.}
\sol{126}{D}{Construyendo una aplicación continua sobre el simplex de precios que sube precios con exceso de demanda, el Teorema de Brouwer (o Kakutani si hay correspondencias) garantiza un punto fijo $p^*$ con $z(p^*) \le 0$.}
\sol{127}{B}{SMD demuestra que la racionalidad individual no impone restricciones suficientes sobre $z(p)$ agregado (salvo continuidad, homogeneidad de grado 0 y Ley de Walras), por lo que puede haber múltiples equilibrios e inestabilidad.}
\sol{128}{C}{La propiedad de sustitubilidad bruta (gross substitutes) implica que si sube $p_j$, aumenta el exceso de demanda de los demás bienes, lo que asegura unicidad y estabilidad global del tâtonnement.}
\sol{129}{A}{El cumplimiento de WARP a nivel agregado (por ejemplo, cuando existe un consumidor representativo positivo) impide que haya múltiples precios relativos de equilibrio aislados.}
\sol{130}{D}{En economías regulares la suma de los índices de los equilibrios es $+1$ (siendo cada índice $+1$ o $-1$), lo que obliga a que exista un número finito e impar de equilibrios.}
\sol{131}{B}{Si $z_1'(p_1^*) < 0$, un precio ligeramente superior a $p_1^*$ genera exceso de oferta ($z_1 < 0$) y el subastador baja $p_1$ de vuelta hacia $p_1^*$.}
\sol{132}{C}{Cuando un agente es oferente neto de un bien, un alza de $p_1$ eleva el valor de su dotación; si ese efecto riqueza sobre su propia demanda domina al efecto sustitución, $dz_1/dp_1$ puede volverse positivo.}
\sol{133}{A}{Arrow-Debreu define una mercancía por sus características físicas, fecha, lugar y estado de la naturaleza, extendiendo el equilibrio general al análisis intertemporal y bajo incertidumbre.}
\sol{134}{D}{Si las preferencias son estrictamente monótonas, cuando $p_j \to 0$ la demanda por el bien $j$ diverge ($\|z(p)\| \to \infty$), asegurando que el equilibrio no quede en la frontera con precios nulos.}
\sol{135}{B}{Stolper-Samuelson (efecto magnificación de Jones) prueba que $\hat{w} > \hat{p}_X > \hat{p}_Y > \hat{r}$: el factor usado intensivamente en el bien que se encarece gana en términos reales de ambos bienes y el otro pierde.}
\sol{136}{C}{Por Rybczynski, a precios constantes $(w/r$ y proporciones $K/L$ fijas), absorber el trabajo extra exige expandir $X$ y retirar capital del sector $Y$, reduciendo absolutamente la producción de $Y$.}
\sol{137}{A}{Con tecnologías idénticas, RCE y sin reversión de intensidades factoriales, la igualdad de precios de bienes $(p_X, p_Y)$ determina unívocamente iguales $(w, r)$ en ambos países.}
\sol{138}{D}{La abundancia relativa de capital abarata $r/w$ en autarquía, dando ventaja comparativa en el bien intensivo en capital.}
\sol{139}{B}{Harberger mostró que gravar el capital en un sector induce sustitución hacia trabajo en $X$ (bajando $r$) y encarece $X$ contrayendo su producción (efecto producto que también presiona $r$ si $X$ es intensivo en $K$).}
\sol{140}{C}{Si $X$ es intensivo en trabajo, la contracción de $X$ (efecto producto) libera relativamente mucho trabajo y poco capital, lo que puede bajar $w/r$ y trasladar la carga al trabajo.}
\sol{141}{A}{Salvo en el caso especial de FPP lineal (costos constantes), el equilibrio general determina $p_X/p_Y$ conjuntamente por la oferta relativa (FPP) y la demanda relativa.}
\sol{142}{D}{Cada punto sobre la FPU corresponde a una asignación Pareto eficiente (Curva de Contrato y eficiencia global); ningún punto fuera de la FPU es factible.}
\sol{143}{B}{Cuando dos FPU asociadas a asignaciones distintas se cruzan, $B$ puede dominar a $A$ según Kaldor-Hicks y simultáneamente $A$ dominar a $B$, generando intransitividad.}
\sol{144}{C}{Para $\rho = 1$ es Bentham (cero aversión a desigualdad en utilidades); para $\rho \to 0$ es Nash ($U_1 U_2$); para $\rho \to -\infty$ es Rawls ($\min\{U_1, U_2\}$, aversión infinita).}
\sol{145}{A}{Arrow probó que cualquier regla transitiva y completa que cumpla dominio universal, Pareto e IIA es necesariamente dictatorial.}
\sol{146}{D}{IIA prohíbe que la presencia o posición relativa de una tercera alternativa $z$ invierta la clasificación social entre $x$ e $y$ (reglas como el conteo de Borda violan IIA).}
\sol{147}{B}{El Teorema de Black (votante mediano) muestra que si las preferencias son unimodales en un espacio unidimensional (y hay número impar de votantes), el pico del votante mediano vence a cualquier otra alternativa.}
\sol{148}{C}{Maximizar $W(U_A, U_B)$ sujeto a la FPU requiere tangencia (en solución interior): el cociente de ponderaciones marginales $(\partial W/\partial U_A)/(\partial W/\partial U_B)$ iguala la pendiente de la FPU.}
\sol{149}{A}{El Segundo Teorema sugiere mantener las señales de precio eficientes ($P = CMg$) y abordar la equidad vía transferencias directas de suma alzada (bonos focalizados), evitando distorsiones por subsidios cruzados excesivos.}
\sol{150}{D}{El Análisis Costo-Beneficio estándar aprueba proyectos con beneficio neto positivo ($\Delta W > 0$), lo que corresponde a una mejora potencial de Pareto (criterio de Kaldor-Hicks).}
\sol{151}{B}{La función de bienestar de Rawls ($W = \min_i U_i$) asigna prioridad absoluta a mejorar la situación de los individuos con menor nivel de bienestar, como los hogares sin acceso básico al agua.}
\sol{152}{C}{Sin mercados completos de bienes contingentes (Arrow-Debreu) para estados secos y húmedos, los agentes no pueden transferir eficientemente el riesgo hidrológico.}
\sol{153}{A}{Las economías de escala en redes (monopolio natural) y las externalidades positivas del saneamiento (salud pública) o negativas de vertimientos violan los supuestos del Primer Teorema.}
\sol{154}{A}{Por la Ley de Walras, $p_1 z_1 + p_2 z_2 + p_3 z_3 = 0 \implies 2(10) + 3(-5) + 5z_3 = 0 \implies 5 + 5z_3 = 0 \implies z_3 = -1$.}
\sol{155}{B}{Igualando $z_1 = 0$: $6(1) - 2p_1^* = 0 \implies p_1^* = 3$. Al reemplazar $(\lambda p_1, \lambda p_2)$, el factor $\lambda$ se cancela, confirmando homogeneidad de grado $0$.}
\sol{156}{C}{Igualando $TMS_A = TMS_B$: $\frac{1}{2}\frac{y_A}{x_A} = 2\frac{12 - y_A}{12 - x_A}$. Con $x_A = 4$ ($x_B = 8$): $\frac{y_A}{8} = \frac{12 - y_A}{4} \implies y_A = 24 - 2y_A \implies y_A = 8$.}
\sol{157}{D}{Con $p_y=1$, $I_A = 12p_x$ y $I_B = 6$, dando $y_A = 6p_x$ y $y_B = 3$. Vaciando el mercado de $y$: $6p_x + 3 = 6 \implies p_x^* = 0.5$, $x_A^* = 6$ y $y_A^* = 3$.}
\sol{158}{A}{Con $p_y=1$, $x_A = \frac{1}{4}\frac{4p_x+4}{p_x} = 1 + \frac{1}{p_x}$ y $x_B = \frac{3}{4}\frac{4p_x+12}{p_x} = 3 + \frac{9}{p_x}$. Igualando $x_A + x_B = 8$: $4 + \frac{10}{p_x} = 8 \implies p_x = 2.50$.}
\sol{159}{B}{De $TMS_i = p_x$ se obtiene $x_A = 6/p_x$ y $x_B = 10/p_x$. Vaciando el mercado de $x$: $16/p_x = 8 \implies p_x^* = 2$, resultando $x_A^* = 3$ y $x_B^* = 5$ (con $y_A^*=8, y_B^*=12$).}
\sol{160}{C}{En la dotación inicial $U_A(W) = U_B(W) = 9$. Sobre la curva de contrato ($y_A = x_A$), la racionalidad individual exige $x_A^2 \ge 9 \implies x_A \ge 3$ y $(10-x_A)^2 \ge 9 \implies x_A \le 7$.}
\sol{161}{D}{En $A=(7,7)$, $TMS_A = TMS_B = 1$, luego $p_x^* = p_y = 1$. A esos precios, el ingreso inicial de $A$ es $4+2=6$ y su canasta objetivo cuesta $7+7=14$, requiriendo una transferencia $T = 14 - 6 = 8$.}
\sol{162}{A}{Como la dotación total de $x$ ($10$) excede a la de $y$ ($6$) y ambos requieren proporción $1:1$, el bien $x$ es abundante (bien libre), luego $p_x^* = 0$. Así $I_A = 4$, $y_A^* = 4$ y $U_A^* = 4$.}
\sol{163}{B}{La eficiencia técnica exige $TMST_X = TMST_Y$: $\frac{K_X}{L_X} = \frac{1}{2}\frac{90 - K_X}{90 - L_X}$. Con $L_X = 45$ ($L_Y = 45$): $K_X = \frac{1}{2}(90 - K_X) \implies 1.5K_X = 45 \implies K_X = 30$.}
\sol{164}{C}{Diferenciando la FPP: $2X dX + 8Y dY = 0 \implies TMT = X/(4Y)$. Igualando a $TMS = Y/X$ resulta $X^2 = 4Y^2$. Sustituyendo en la FPP: $8Y^2 = 200 \implies Y^*=5$, $X^*=10$ y $p_X/p_Y = 5/10 = 0.5$.}
\sol{165}{D}{La FPP es $X/2 + Y/4 = 30 \iff 2X + Y = 120$, por lo que $TMT = p_X/p_Y = 2$. Igualando $TMS = Y/X = 2$ se tiene $Y = 2X$, dando $4X = 120 \implies X^* = 30$ y $Y^* = 60$.}
\sol{166}{A}{En $(4, 8)$: $TMT = CMg_X/CMg_Y = 8/32 = 0.25$ y $TMS = 8/4 = 2.0$. Como $TMS > TMT$, la valoración marginal de $X$ supera su costo de oportunidad: conviene producir más $X$ y menos $Y$.}
\sol{167}{B}{Igualando $TMS_{l,c} = PMg_L$: $\frac{c}{l} = \frac{3}{\sqrt{L}}$. Sustituyendo $c = 6\sqrt{L}$ y $l = 27 - L$: $\frac{6\sqrt{L}}{27 - L} = \frac{3}{\sqrt{L}} \implies 2L = 27 - L \implies L^* = 9$ y $c^* = 6(3) = 18$.}
\sol{168}{C}{Sobre la FPP, $TMT = p_X^p/p_Y = 1/2$. Con impuesto del $100\%$, $p_X^c/p_Y = 1$, por lo que $TMS = Y/X = 1 \implies Y = X$. Sustituyendo en $X + 2Y = 120$ resulta $X_t = 40$ e $Y_t = 40$.}
\sol{169}{D}{Resolviendo $2w_1 + r_1 = 13$ y $w_1 + 2r_1 = 11$: multiplicando la primera por $2$ y restando la segunda da $3w_1 = 15 \implies w_1 = 5$ y $r_1 = 3$ (frente a $w_0=3, r_0=4$).}
\sol{170}{A}{Restando las ecuaciones en variaciones: $2\Delta X + \Delta Y = 6$ y $\Delta X + 2\Delta Y = 0 \implies \Delta X = -2\Delta Y \implies -3\Delta Y = 6 \implies \Delta Y = -2$ y $\Delta X = +4$.}
\sol{171}{B}{La tangencia entre la recta de isobienestar y la FPU exige $\frac{3}{1} = \frac{U_A}{U_B} \implies U_A = 3U_B$. Sustituyendo en la FPU: $10U_B^2 = 160 \implies U_B^* = 4$ y $U_A^* = 12$.}
\sol{172}{C}{Igualando la pendiente de $W$ ($U_B/U_A$) a la pendiente de la FPU ($2$): $U_B = 2U_A$. Sustituyendo en $2U_A + U_B = 120$ da $4U_A = 120 \implies U_A^* = 30$ y $U_B^* = 60$.}
\sol{173}{D}{El óptimo rawlsiano se alcanza en el vértice $U_A = U_B$ sobre la FPU: $U_A + 3U_A = 120 \implies 4U_A = 120 \implies U_A^* = U_B^* = 30$.}
\sol{174}{A}{Igualando la tasa marginal de sustitución social $(U_B/U_A)^{1/2}$ a la pendiente de la FPU ($1/4$): $U_B/U_A = 1/16 \implies U_A = 16U_B$. En la FPU: $20U_B = 100 \implies U_B^* = 5$ y $U_A^* = 80$.}
\sol{175}{B}{Para que ambos grupos queden estrictamente mejor, $B$ debe recibir más que su pérdida ($T > 110$) y $A$ debe pagar menos que su ganancia ($T < 180$), es decir, $110 < T < 180$.}
\sol{176}{C}{Por equimarginalidad, $100 - 2q_1 = 85 - (60 - q_1) \implies 100 - 2q_1 = 25 + q_1 \implies 3q_1 = 75 \implies q_1^* = 25$, $q_2^* = 35$ y $\lambda^* = 100 - 50 = 50$.}
\sol{177}{D}{En $(45, 15)$, $BMg_1 = 10$ y $BMg_2 = 70$. Al transferir $20$ unidades de $1$ a $2$ hasta igualar ambos en $50$, el triángulo de ganancia neta es $\Delta W = \frac{1}{2}(70 - 10)(20) = 600$.}
\sol{178}{A}{Por el Teorema del Votante Mediano (preferencias unimodales con número impar de votantes), el ganador de Condorcet es la mediana de los picos ($S/\ 50$), no el promedio ($S/\ 61$).}
\sol{179}{B}{Como ambos tienen preferencias idénticas homotéticas, en equilibrio $c_{iH}/c_{iS} = p_S/p_H = C_H/C_S = (12+4)/(4+4) = 16/8 = 2.0$ (el agua es dos veces más cara en sequía).}
\sol{180}{C}{Con $p_x'=4$, el hogar elige $x'=9, y'=36$ ($U'=18$) y el subsidio cuesta $(9-4)(9) = 45$. A precios $(9,1)$, el gasto mínimo para $U'=18$ es $e = 2\sqrt{9\cdot 1}(18) = 108$, requiriendo una suma alzada de $108 - 72 = 36$ (ahorro de $9$).}
